\documentclass[11pt]{article}

% ---------- Packages ----------
\usepackage[top=0.7in,bottom=0.7in,left=0.75in,right=0.75in]{geometry}
\usepackage{titlesec}
\usepackage{xcolor}
\usepackage{tabularx}
\usepackage{enumitem}
\usepackage{fontawesome5}
\usepackage{hyperref}
\usepackage{charter}
\usepackage{longtable}

\pagenumbering{gobble}
\setlength{\parindent}{0pt}
\setlist[itemize]{noitemsep,topsep=2pt}

% ---------- Colors ----------
\definecolor{headingblue}{RGB}{60,110,170}
\definecolor{namegray}{RGB}{0,0,0}

% Keep body text + links black
\hypersetup{colorlinks=true,linkcolor=black,urlcolor=black,citecolor=black}

% ---------- Section Style ----------
\newcommand{\cvsection}[1]{%
\vspace{12pt}%
\noindent%
{\color{headingblue}\rule{40pt}{5pt}}%
\hspace{10pt}%
{\Large\textbf{\color{headingblue}#1}}%
\vspace{6pt}%
\color{black}%
}

% =====================================================
\begin{document}

% =====================================================
% HEADER
% =====================================================

\begin{tabularx}{\textwidth}{Xr}

{\fontsize{36}{36}\selectfont\color{namegray} DEVESH KAUSHAL}

&
\small
\faPhone\ 9818535128 \\
&
\faEnvelope\ Deveshkaushal9@gmail.com \\
&
\faLinkedin\ LinkedIn \\
&
\faGraduationCap\ Google Scholar \\

\end{tabularx}

\vspace{6pt}
% {\Large\itshape Curriculum Vitae}

% =====================================================
% EDUCATION
% =====================================================
% =====================================================
% RESEARCH INTERESTS
% =====================================================

\cvsection{Research Interests}

AI Systems Assurance and Evaluation, Safety-Critical Machine Learning, Robustness and Failure Modes of Autonomous Systems, Multi-Agent System Evaluation, Physical-World AI Modeling, AI for Real-World Risk Prediction
\cvsection{Education}

\begin{tabularx}{\textwidth}{p{90pt}X}

2022--2026 &
\textbf{Bachelor of Technology, Computer Science and Engineering} \\
& I.K. Gujral Punjab Technical University \\

\end{tabularx}

% =====================================================
% RESEARCH EXPERIENCE
% =====================================================

\cvsection{Research Experience}

\begingroup
\renewcommand{\arraystretch}{1.1}
\begin{tabularx}{\textwidth}{@{}p{95pt}X@{}}

{\small\color{namegray}Present} &
\textbf{Indian Institute of Technology Delhi}\hfill\textit{Undergraduate Research Assistant}

\begin{itemize}[leftmargin=*,nosep]
\item Working on FALLSAFE drone and AI terrain risk mapping framework.
\item Developing UAV terrain modelling pipelines producing DSM, DOM, and 3D surface reconstructions.
\item Integrating geological modelling and GIS workflows for slope instability analysis.
\item Performing failed slope back-analysis using drone reconstructions.
\item Supporting ML-based terrain risk prediction models.
\item Applied machine learning models to support data-driven risk assessment and operational decision-making in real-world settings.

\end{itemize}
\\[6pt]

{\small\color{namegray}June--Aug 2025} &
\textbf{Indian Institute of Technology Ropar}\hfill\textit{Research Intern}

\begin{itemize}[leftmargin=*,nosep]
\item Conducted wireless sensing research using WiFi CSI.
\item Developed LightPoseNet pose estimation model.
\item Paper accepted at IEEE Massachusetts Institute of Technology URTC 2025.
\item Built CSI preprocessing pipelines.
\end{itemize}
\\[6pt]

{\small\color{namegray}July--Sept 2025} &
\textbf{Centre for National Security Studies}\hfill\textit{Research Intern}

\begin{itemize}[leftmargin=*,nosep]
\item Studied emerging defence technologies and UAV surveillance.
\item Authored research articles on national security innovations.
\end{itemize}

\end{tabularx}
\endgroup



% =====================================================
% PUBLICATIONS & PRESENTATIONS
% =====================================================

% =====================================================
% CONFERENCE PRESENTATIONS & PUBLICATIONS
% =====================================================

\cvsection{Conference Presentations \& Publications}

\begingroup
\renewcommand{\arraystretch}{1.1}
% \begin{tabularx}{\textwidth}{@{}p{95pt}X@{}}
\begin{longtable}{@{}p{95pt}p{\dimexpr\textwidth-95pt\relax}@{}}

Jan 2026 &
\textbf{AAAI – Empowering Global South AI (EGSAI Workshop)} \hfill \textit{Singapore}

\begin{itemize}[leftmargin=*,nosep]
\item Poster accepted at AAAI 2026 EGSAI Workshop.
\item Invited to submit camera-ready version and 3-minute research presentation.
\end{itemize}
\\[6pt]

Dec 2025 &
\textbf{IEEE International Conference on Edge Computing (ICEdge)} \hfill \textit{IISc Bangalore}

\begin{itemize}[leftmargin=*,nosep]
\item Accepted for R\&D Showcase and Student Forum.
\item Presented \textit{LightPoseNet: Lightweight Human Pose Estimation from WiFi Signals}.
\item Demonstrated privacy-preserving 3D pose estimation using commodity WiFi.
\end{itemize}
\\[6pt]

Dec 2025 &
\textbf{TechConnect – ResCon Poster Competition} \hfill \textit{IIT Bombay}

\begin{itemize}[leftmargin=*,nosep]
\item Selected finalist for poster presentation round.
\item Presented WiFi-based privacy-preserving 3D pose estimation research.
\end{itemize}
\\[6pt]

Oct 2025 &
\textbf{IEEE MIT Undergraduate Research Technology Conference (URTC)} \hfill \textit{Cambridge, MA}

\begin{itemize}[leftmargin=*,nosep]
\item Accepted paper: \textit{LightPoseNet: Lightweight Human Pose Estimation from WiFi CSI}.
\item Developed real-time WiFi-based 3D pose estimation evaluated on WiPose dataset.
\end{itemize}
\\[6pt]

June 2025 &
\textbf{International Conference on Electronics, AI and Computing (EAIC)}\hfill \textit{NIT Jalandhar}

\begin{itemize}[leftmargin=*,nosep]
\item Paper: \textit{Radar-Based UAV Classification: A Micro-Doppler and Deep Learning Approach}.
\item Proposed deep learning model for UAV classification using radar micro-Doppler signatures.
\end{itemize}
\\[6pt]

May 2025 &
\textbf{Interdisciplinary Conference on AI \& Cognitive Science} \hfill \textit{University of Cambridge, UK}

\begin{itemize}[leftmargin=*,nosep]
\item Presented airborne object detection using EfficientNetB3.
\item Selected as one of 65 global presenters.
\end{itemize}
\\[6pt]

March 2025 &
\textbf{International Conference on Statistical Learning (ICSL)} \hfill \textit{IIT Ropar}

\begin{itemize}[leftmargin=*,nosep]
\item Presented signal processing and ML methods for UAV detection.
\end{itemize}

\end{longtable}
\endgroup

% =====================================================
% RESEARCH PROJECTS
% =====================================================
% =====================================================
% RESEARCH PROJECTS
% =====================================================

\cvsection{Research Projects}

\begingroup
\renewcommand{\arraystretch}{1.1}
\begin{tabularx}{\textwidth}{@{}p{95pt}X@{}}

2026 &
\textbf{Geometry-Aware Terrain Reconstruction using Drone Photogrammetry and Learning-Based Surface Modeling}

\begin{itemize}[leftmargin=*,nosep]
\item Developing a pipeline for reconstructing high-resolution 3D terrain surfaces from UAV photogrammetry data.
\item Integrating classical geometry processing techniques with deep learning for surface smoothing and feature extraction.
\item Exploring mesh representations, point-cloud processing, and geometric feature descriptors for terrain risk analysis.
\item Applying computational geometry methods to enhance geospatial modeling and data-driven slope instability assessment.
\item Studying discrete differential geometry concepts including curvature estimation, Laplacian smoothing, and mesh optimization.
\item \textbf{Tools:} Open3D, PyVista, PyMeshLab; self-studied discrete differential geometry (curvature, Laplace–Beltrami operators, mesh processing) using Keenan Crane’s CMU course materials.

\end{itemize}
\\[6pt]

2025 &
\textbf{LightPoseNet: 3D Human Pose Estimation from WiFi CSI}

\begin{itemize}[leftmargin=*,nosep]
\item Developed lightweight CNN framework for WiFi-based 3D human pose estimation.
\item Proposed baseline and SE-attention variants achieving \textbf{40 mm MPJPE} and \textbf{93.31\% PCK@0.10}.
\item Achieved \textbf{300× parameter reduction} enabling real-time edge deployment.
\item Designed privacy-preserving, through-wall pose sensing using commodity WiFi.
\end{itemize}
\\[6pt]

2025 &
\textbf{Radar-Based UAV Classification: Micro-Doppler and Deep Learning}

\begin{itemize}[leftmargin=*,nosep]
\item Developed deep learning pipeline for radar-based drone vs bird classification.
\item Fine-tuned \textbf{EfficientNetB3} achieving \textbf{97.10\% classification accuracy}.
\item Utilized DIAT-µSAT micro-Doppler dataset from Defence Institute of Advanced Technology.
\end{itemize}
\\[6pt]

2024 &
\textbf{AI-Powered Disaster Prediction Using Geospatial Data}

\begin{itemize}[leftmargin=*,nosep]
\item Presented system at IIT Kanpur disaster technology showcase.
\item Built ML pipeline combining satellite imagery and environmental data.
\item Integrated GIS features for spatio-temporal flood and landslide forecasting.
\end{itemize}
\\[6pt]



\end{tabularx}
\endgroup

% =====================================================
% TECHNICAL SKILLS
% =====================================================

% =====================================================
% TECHNICAL SKILLS
% =====================================================

\cvsection{Technical Skills}

\textbf{Programming:} Python, Dart, SQL, JavaScript, C, Java \\
\textbf{Machine Learning \& AI:} PyTorch, TensorFlow, scikit-learn, OpenCV, HuggingFace Transformers, LangChain, LLaMA, OpenAI API, Pinecone \\
\textbf{Developer Tools:} Git, GitHub, Docker, Firebase, Jupyter Notebooks, Android Studio, Postman, VS Code



% =====================================================
% SUMMER SCHOOLS & ADVANCED PROGRAMS
% =====================================================
\newpage
\cvsection{Fellowship}

\begingroup
\renewcommand{\arraystretch}{1.05}
\begin{longtable}{@{}p{95pt}p{\dimexpr\textwidth-95pt\relax}@{}}

2025 &
\textbf{ACM India COMPUTE Fellow} \hfill \textit{IIT Ropar}
\begin{itemize}[leftmargin=*,nosep]
\item Selected COMPUTE 2025 Fellow for ACM India and iSIGCSE computing education conference.
\item Participated in technical sessions, expert talks, and academic networking events.
\end{itemize}
\\[4pt]

% \end{itemize}

\end{longtable}
\endgroup
\cvsection{Summer Schools \& Advanced Programs}

\begingroup
\renewcommand{\arraystretch}{1.05}
\begin{longtable}{@{}p{95pt}p{\dimexpr\textwidth-95pt\relax}@{}}



2025 &
\textbf{EE Summer School} \hfill \textit{Indian Institute of Science (IISc), Bangalore}
\begin{itemize}[leftmargin=*,nosep]
\item Selected among 100 students nationwide.
\item Completed advanced modules in signal processing, control systems, power electronics, and renewable energy integration.
\end{itemize}
\\[4pt]

2025 &
\textbf{IEEE–EURASIP Summer School on Signal Processing (S3P)} \hfill \textit{Tuscany, Italy}
\begin{itemize}[leftmargin=*,nosep]
\item Selected undergraduate participant for advanced research-focused signal processing program designed primarily for PhD researchers.
\end{itemize}
\\[4pt]

2025 &
\textbf{Oxford Machine Learning Summer School (OxML)} \hfill \textit{University of Oxford}
\begin{itemize}[leftmargin=*,nosep]
\item Accepted for intensive coursework in machine learning, deep learning, and AI applications.
\item Selected among competitive global applicants.
\end{itemize}

\end{longtable}
\endgroup


\cvsection{Referees}

\begin{minipage}[t]{0.45\textwidth}
\textbf{Dr. Sumeet Kumar Sinha}\\
\textit{Assistant Professor}\\
Dept. of Civil \& Environmental Engineering\\
IIT Delhi\\
\faEnvelope\ sksinha@civil.iitd.ac.in
\end{minipage}
\hfill
\begin{minipage}[t]{0.45\textwidth}
\textbf{Prof. (Dr.) Vikas Chawla}\\
\textit{Director}\\
I.K. Gujral Punjab Technical University\\
\faEnvelope\ vikas.chawla@ptu.ac.in
\end{minipage}
\hfill




\cvsection{Grants \& Funding}

\begingroup
\renewcommand{\arraystretch}{1.1}
\begin{tabularx}{\textwidth}{@{}p{95pt}X@{}}

2026 &
\textbf{Travel Grant – TreeHacks} \hfill \textit{Stanford University}

\begin{itemize}[leftmargin=*,nosep]
\item Awarded travel reimbursement grant to participate in TreeHacks 2026.
\item Selected for highly competitive international hackathon hosted at Stanford University.
\end{itemize}
\\[6pt]

2025 &
\textbf{Travel Grant – Code for Good Hackathon} \hfill \textit{IIT Mandi}

\begin{itemize}[leftmargin=*,nosep]
\item Received travel and accommodation grant from I.K. Gujral Punjab Technical University.
\item Selected among top 10 teams nationwide based on technical innovation.
\end{itemize}
\\[6pt]

2025 &
\textbf{IISc EE Summer School Grant} \hfill \textit{Indian Institute of Science, Bangalore}

\begin{itemize}[leftmargin=*,nosep]
\item Awarded competitive grant to attend IISc EE Summer School 2025.
\item Selected among 100 participants nationwide.
\item Completed advanced modules in signal processing, control systems and power electronics.
\end{itemize}

\end{tabularx}
\endgroup


% =====================================================
% LEADERSHIP
% =====================================================

% =====================================================
% LEADERSHIP & EXTRACURRICULAR
% =====================================================
\newpage
\cvsection{Leadership \& Extracurricular}

\begingroup
\renewcommand{\arraystretch}{1.1}
\begin{tabularx}{\textwidth}{@{}p{95pt}X@{}}

2023--2024 &
\textbf{Google Developer Student Clubs (GDSC)} \hfill \textit{Android Core Team Member, IKGPTU Main Campus}

\begin{itemize}[leftmargin=*,nosep]
\item Organized workshops and technical events on Android development.
\item Conducted peer training sessions on mobile application deployment.
\item Mentored students during hackathons and coding competitions.
\end{itemize}
\\[6pt]

2024--2025 &
\textbf{Bizarre Coder} \hfill \textit{Team Lead, IKGPTU Main Campus}

\begin{itemize}[leftmargin=*,nosep]
\item Coordinated technical teams and organized coding workshops.
\item Mentored students participating in hackathons and collaborative projects.
\item Facilitated project-based coding sessions to enhance real-world development skills.
\end{itemize}

\end{tabularx}
\endgroup

% =====================================================
% ACHIEVEMENTS
% =====================================================
% =====================================================
% ACHIEVEMENTS
% =====================================================

\cvsection{Achievements}

\begin{itemize}[leftmargin=*,nosep]
\item Selected participant — \textbf{MIT Hacking Medicine Grand Hack 2026} (\textbf{MIT}, Boston);wil contribute to the \textbf{Trauma \& Rehabilitation} track, developing rapid, interdisciplinary healthcare innovation solutions.

\item Selected for \textbf{TreeHacks 2026} at \textbf{Stanford University}, chosen from 15,000+ global applicants.
\item{Selected to participate in the \textbf{MIT iQuHACK 2026 Remote Competition}, one of the world’s largest student-run \textbf{quantum computing hackathons}, hosted by \textbf{Massachusetts Institute of Technology (MIT)}.}

\item Selected presenter at the \textbf{Interdisciplinary Conference on AI \& Cognitive Science}, University of Cambridge, UK.
  \item{Selected as a \textbf{COMPUTE 2025 Fellow} for the 8\textsuperscript{th} ACM India and iSIGCSE conference on computing education, held at \textbf{IIT Ropar}.}

  \item {Selected for the \textbf{Online Winter Design Camp} organized by the \textbf{Design Experience Lab, IIT Ropar}, powered by \textbf{Dassault Syst\`emes}.}

\item \textbf{1st Place} – Internal Smart India Hackathon (SIH) selection round.

\item \textbf{Finalist} – Code for Good Hackathon, IIT Mandi (Top 10 of 550 teams).

\item \textbf{Finalist} – Hackfest, IIT (ISM) Dhanbad (Top 70 of 1,200 teams).

\item \textbf{Top 40 Teams} – IIIT Naya Raipur Hackathon (1,764 teams).

\item \textbf{Finalist} – Hack’Ndore (Top 100 of 1,100 teams).

\item \textbf{Top 400 Teams} – HackWithIndia (1,500 teams).

\item Advanced to Round 2 – Haryana Police Narcotics Hackathon, IGDTUW Hackathon, and Bank of Baroda Hackathon.

\item Mentored 50+ students in mobile app development and machine learning during a national hackathon (Final round hosted at Microsoft Office, Gurgaon).

\end{itemize}

% =====================================================
\end{document}
